\documentclass[11pt,a4paper]{article}

% ── Codificación, fuentes y lenguaje ───────────────────────────────────────────
\usepackage[utf8]{inputenc}
\usepackage[T1]{fontenc}
\usepackage[default,scale=0.95]{sourcesanspro}
\renewcommand{\familydefault}{\sfdefault}
\usepackage[spanish,es-noshorthands,es-tabla]{babel}

% ── Geometría y espaciado ─────────────────────────────────────────────────────
\usepackage[top=2.2cm, bottom=2.5cm, left=2.2cm, right=2.2cm, headheight=15pt]{geometry}
\usepackage{setspace}
\setstretch{1.18}
\usepackage{parskip}

% ── Matemáticas y unidades ────────────────────────────────────────────────────
\usepackage{amsmath}
\usepackage{amssymb}
\usepackage{siunitx}

% ── Circuitos ─────────────────────────────────────────────────────────────────
\usepackage[european, straightvoltages]{circuitikz}

% ── Tablas ────────────────────────────────────────────────────────────────────
\usepackage{booktabs}
\usepackage{tabularx}
\usepackage{array}
\usepackage{longtable}
\usepackage{needspace}
\usepackage{colortbl}
\usepackage{enumitem}

% ── Colores, cajas e iconos ───────────────────────────────────────────────────
\usepackage[dvipsnames]{xcolor}
\usepackage{tcolorbox}
\tcbuselibrary{skins, breakable}
\usepackage{fontawesome5}
\usepackage{graphicx}

% ── Listados de código (dark-mode infográfico) ────────────────────────────────
%   Estilo base: fondo oscuro con números de línea. Los bloques lstlisting se
%   envuelven automáticamente en una caja tcolorbox redondeada.
\usepackage{listings}

% ── Secciones con barra de color (infografía) ────────────────────────────────
\usepackage{titlesec}
\titleformat{\section}
  {\normalfont\LARGE\bfseries\color{azulNoche}}
  {\colorbox{cyanNeon}{\color{fondoOscuro}\thesection}}{0.7em}{}
  [\vspace{2pt}\color{cyanNeon}\rule{\linewidth}{1.8pt}]
\titlespacing*{\section}{0pt}{24pt}{10pt}
\titleformat{\subsection}
  {\normalfont\large\bfseries\color{azulNoche}}
  {\textcolor{cyanNeon}{\thesubsection}}{0.7em}{}
  [\vspace{1pt}\color{grisLinea}\rule{\linewidth}{0.6pt}]
\titlespacing*{\subsection}{0pt}{14pt}{6pt}
\titleformat{\subsubsection}
  {\normalfont\normalsize\bfseries\color{azulNoche}}
  {\textcolor{cyanNeon}{\thesubsubsection}}{0.7em}{}
\titlespacing*{\subsubsection}{0pt}{10pt}{4pt}

% ── Cabeceras y pies de página ────────────────────────────────────────────────
\usepackage{fancyhdr}
\usepackage{lastpage}
\pagestyle{fancy}
\fancyhf{}
\renewcommand{\headrulewidth}{0pt}
\renewcommand{\footrulewidth}{0pt}
\fancyfoot[C]{%
  \begin{minipage}{\linewidth}
    \vspace{2pt}
    \color{cyanNeon}\rule{\linewidth}{1.2pt}\\[2pt]
    \centering\color{grisTexto}\footnotesize
    Página \thepage\ de \pageref{LastPage}
  \end{minipage}%
}

% ── Hiperenlaces ──────────────────────────────────────────────────────────────
\usepackage[colorlinks=true, linkcolor=azulNoche, urlcolor=cyanNeon]{hyperref}
\sloppy
\emergencystretch 3em

% ── Paleta de colores — tecnológico dark-mode ──────────────────────────────────
\definecolor{fondoOscuro}{HTML}{0D1B2A}%             Dark navy — fondo banda título
\definecolor{verdeTurquesa}{HTML}{004D40}%           Verde turquesa — bandas de marca
\definecolor{azulNoche}{HTML}{1B3A6B}%               Midnight blue — secciones, marcos
\definecolor{azulMedio}{HTML}{2A5298}%               Azul medio — variante de acento
\definecolor{cyanNeon}{HTML}{00C8E0}%                Cyan eléctrico — acentos primarios
\definecolor{verdeSignal}{HTML}{00B887}%             Verde señal — fortalezas, éxito
\definecolor{naranjaVivo}{HTML}{FF7043}%             Naranja vivo — mejoras, advertencias
\definecolor{rojoAlerta}{HTML}{E53935}%              Rojo alerta — errores
\definecolor{grisPapel}{HTML}{F4F6FA}%               Fondo tarjetas claras
\definecolor{grisLinea}{HTML}{C8D0E0}%               Bordes sutiles
\definecolor{grisTexto}{HTML}{6B7A99}%               Texto secundario
\definecolor{amarilloNota}{HTML}{FFB300}%            Notas / advertencias doradas
\definecolor{codigoFondo}{HTML}{1E2D3D}%             Fondo del bloque de código
\definecolor{codigoComentario}{HTML}{7A8FA6}%        Comentarios
\definecolor{codigoCadena}{HTML}{FF9D5C}%            Cadenas / strings
\definecolor{codigoKeyword}{HTML}{00C8E0}%           Palabras clave
\definecolor{codigoNumero}{HTML}{00E5A0}%            Números

% ── Banda de título: portada infográfica ──────────────────────────────────────
%   Diseño en 2 capas: banda de color (por defecto fondo oscuro) + franja de
%   acento cyan abajo. Uso: \bandaTitulo[color]{título}{subtítulo}.
%   El icono se puede cambiar con \renewcommand{\iconoBanda}{...} (FontAwesome).
\newcommand{\iconoBanda}{microchip}
\newcommand{\bandaTitulo}[3][fondoOscuro]{%
  \begin{tcolorbox}[
    enhanced,
    colback=#1, colframe=#1,
    boxrule=0pt, arc=8pt, outer arc=8pt,
    width=\linewidth,
    top=18pt, bottom=6pt, left=14pt, right=14pt,
    borderline south={3.5pt}{0pt}{cyanNeon},
  ]
    \begin{minipage}[c]{0.07\linewidth}
      \centering
      \textcolor{cyanNeon}{\fontsize{30}{30}\selectfont\faIcon{\iconoBanda}}
    \end{minipage}%
    \hspace{8pt}%
    \begin{minipage}[c]{0.88\linewidth}
      {\fontsize{20}{24}\selectfont\bfseries\color{white}#2}\par\vspace{5pt}%
      \textcolor{cyanNeon!80!white}{\small\faIcon{angle-right}~#3}%
    \end{minipage}
    \vspace{4pt}
  \end{tcolorbox}%
  \vspace{8pt}%
}

% ── Tarjetas de datos (infografía) ────────────────────────────────────────────
\tcbset{
  tarjetaDato/.style={
    enhanced, breakable,
    arc=6pt, outer arc=6pt,
    colback=grisPapel, colframe=azulNoche,
    boxrule=1pt,
    fonttitle=\bfseries\small, coltitle=white,
    attach boxed title to top left={yshift=-3mm, xshift=6mm},
    boxed title style={
      colback=azulNoche, colframe=azulNoche,
      arc=4pt, boxrule=0pt,
    },
    drop fuzzy shadow=azulNoche!30!white,
    top=8pt, bottom=8pt, left=10pt, right=10pt,
  },
  tarjetaIcono/.style={
    enhanced, breakable,
    arc=6pt, outer arc=6pt,
    colback=white, colframe=grisLinea, boxrule=0.8pt,
    fonttitle=\bfseries\small, coltitle=azulNoche,
    borderline west={3pt}{0pt}{cyanNeon},
    drop fuzzy shadow=azulNoche!25!white,
    top=6pt, bottom=6pt, left=10pt, right=10pt,
  },
  cajaTitulo/.style={
    enhanced, breakable,
    arc=6pt, outer arc=6pt,
    colback=azulNoche!10!grisPapel, colframe=azulNoche,
    boxrule=1pt,
    fonttitle=\bfseries, coltitle=white,
    attach boxed title to top left={yshift=-3mm, xshift=6mm},
    boxed title style={
      colback=azulNoche, colframe=azulNoche,
      arc=4pt, boxrule=0pt,
    },
    drop fuzzy shadow=azulNoche!25!white,
    top=8pt, bottom=8pt, left=10pt, right=10pt,
  },
  cajaContenido/.style={
    enhanced, breakable,
    arc=5pt, outer arc=5pt,
    colback=grisPapel, colframe=grisLinea,
    boxrule=0.8pt,
    fonttitle=\bfseries\small, coltitle=azulNoche,
    top=6pt, bottom=6pt, left=10pt, right=10pt,
  },
  cajaFortaleza/.style={
    enhanced, breakable,
    arc=6pt, outer arc=6pt,
    colback=verdeSignal!8!white, colframe=verdeSignal,
    boxrule=1pt,
    borderline west={4pt}{0pt}{verdeSignal},
    fonttitle=\bfseries\small, coltitle=verdeSignal!70!black,
    drop fuzzy shadow=verdeSignal!20!white,
    top=6pt, bottom=6pt, left=10pt, right=10pt,
  },
  cajaMejora/.style={
    enhanced, breakable,
    arc=6pt, outer arc=6pt,
    colback=naranjaVivo!8!white, colframe=naranjaVivo,
    boxrule=1pt,
    borderline west={4pt}{0pt}{naranjaVivo},
    fonttitle=\bfseries\small, coltitle=naranjaVivo!80!black,
    drop fuzzy shadow=naranjaVivo!20!white,
    top=6pt, bottom=6pt, left=10pt, right=10pt,
  },
  cajaRecomendacion/.style={
    enhanced, breakable,
    arc=6pt, outer arc=6pt,
    colback=cyanNeon!6!white, colframe=cyanNeon!60!azulNoche,
    boxrule=1pt,
    borderline west={4pt}{0pt}{cyanNeon},
    fonttitle=\bfseries\small, coltitle=azulNoche,
    drop fuzzy shadow=azulNoche!20!white,
    top=6pt, bottom=6pt, left=10pt, right=10pt,
  },
}

% ── Ícono + texto (encabezados de sección y elementos) ────────────────────────
\newcommand{\iconotexto}[2]{\textcolor{cyanNeon}{\faIcon{#1}}~#2}

% ── Listados de código: estilo dark + auto-envoltura ─────────────────────────
\lstdefinestyle{estiloCodigo}{
    backgroundcolor=\color{codigoFondo},
    basicstyle=\ttfamily\small\color{white},
    commentstyle=\color{codigoComentario}\itshape,
    stringstyle=\color{codigoCadena},
    keywordstyle=\color{codigoKeyword}\bfseries,
    numberstyle=\tiny\color{codigoComentario},
    numbers=left,
    numbersep=10pt,
    showstringspaces=false,
    breaklines=true,
    breakatwhitespace=false,
    tabsize=4,
    frame=none,
    captionpos=b,
    aboveskip=6pt,
    belowskip=4pt,
    xleftmargin=14pt,
    framexleftmargin=14pt,
    literate=
      {á}{{\'a}}1 {é}{{\'e}}1 {í}{{\'i}}1 {ó}{{\'o}}1 {ú}{{\'u}}1
      {Á}{{\'A}}1 {É}{{\'E}}1 {Í}{{\'I}}1 {Ó}{{\'O}}1 {Ú}{{\'U}}1
      {ñ}{{\~n}}1 {Ñ}{{\~N}}1 {ü}{{\"u}}1 {Ü}{{\"U}}1
      {¿}{{?`}}1 {¡}{{!`}}1
}
\lstdefinestyle{estiloPython}{
    style=estiloCodigo,
    language=Python,
}
\lstset{style=estiloCodigo}
\BeforeBeginEnvironment{lstlisting}{%
  \begin{tcolorbox}[
    enhanced, breakable,
    arc=6pt, outer arc=6pt,
    colback=codigoFondo, colframe=azulNoche,
    boxrule=1pt,
    drop fuzzy shadow=azulNoche!25!white,
    top=2pt, bottom=2pt, left=0pt, right=0pt,
    colupper=white, coltext=white,
  ]%
}
\AfterEndEnvironment{lstlisting}{\end{tcolorbox}}

% ── Cajas de contenido temáticas ──────────────────────────────────────────────
\newtcolorbox{cajaRecuerda}{
    enhanced, breakable,
    arc=6pt, outer arc=6pt,
    colback=azulNoche!10!grisPapel,
    colframe=azulNoche, boxrule=1pt,
    borderline west={4pt}{0pt}{cyanNeon},
    fonttitle=\bfseries\small\color{white},
    title={\faIcon{info-circle}~Recuerda},
    attach boxed title to top left={yshift=-3mm, xshift=6mm},
    boxed title style={colback=azulNoche, arc=4pt, boxrule=0pt},
    drop fuzzy shadow=azulNoche!25!white,
    left=10pt, right=10pt, top=8pt, bottom=8pt,
}

\newtcolorbox{cajaConcepto}{
    enhanced, breakable,
    arc=6pt, outer arc=6pt,
    colback=verdeSignal!8!white,
    colframe=verdeSignal, boxrule=1pt,
    borderline west={4pt}{0pt}{verdeSignal},
    fonttitle=\bfseries\small\color{white},
    title={\faIcon{lightbulb}~Concepto clave},
    attach boxed title to top left={yshift=-3mm, xshift=6mm},
    boxed title style={colback=verdeSignal!80!black, arc=4pt, boxrule=0pt},
    drop fuzzy shadow=verdeSignal!20!white,
    left=10pt, right=10pt, top=8pt, bottom=8pt,
}

\newtcolorbox{cajaEjemplo}{
    enhanced, breakable,
    arc=6pt, outer arc=6pt,
    colback=cyanNeon!5!white,
    colframe=cyanNeon!70!azulNoche, boxrule=1pt,
    borderline west={4pt}{0pt}{cyanNeon},
    fonttitle=\bfseries\small\color{fondoOscuro},
    title={\faIcon{code}~Ejemplo},
    attach boxed title to top left={yshift=-3mm, xshift=6mm},
    boxed title style={colback=cyanNeon, arc=4pt, boxrule=0pt},
    drop fuzzy shadow=azulNoche!20!white,
    left=10pt, right=10pt, top=8pt, bottom=8pt,
}

\newtcolorbox{cajaEjercicio}{
    enhanced, breakable,
    arc=6pt, outer arc=6pt,
    colback=azulMedio!7!white,
    colframe=azulMedio, boxrule=1pt,
    borderline west={4pt}{0pt}{azulMedio},
    fonttitle=\bfseries\small\color{white},
    title={\faIcon{pencil-alt}~Ejercicio},
    attach boxed title to top left={yshift=-3mm, xshift=6mm},
    boxed title style={colback=azulMedio, arc=4pt, boxrule=0pt},
    drop fuzzy shadow=azulNoche!20!white,
    left=10pt, right=10pt, top=8pt, bottom=8pt,
}

\newtcolorbox{cajaReto}{
    enhanced, breakable,
    arc=6pt, outer arc=6pt,
    colback=naranjaVivo!6!white,
    colframe=naranjaVivo, boxrule=1pt,
    borderline west={4pt}{0pt}{naranjaVivo},
    fonttitle=\bfseries\small\color{white},
    title={\faIcon{fire}~Reto},
    attach boxed title to top left={yshift=-3mm, xshift=6mm},
    boxed title style={colback=naranjaVivo, arc=4pt, boxrule=0pt},
    drop fuzzy shadow=naranjaVivo!20!white,
    left=10pt, right=10pt, top=8pt, bottom=8pt,
}

% ── Indicadores de dificultad ─────────────────────────────────────────────────
\newcommand{\facil}{\textcolor{verdeSignal}{\faIcon{circle}\,\textbf{Fácil}}}
\newcommand{\intermedio}{\textcolor{naranjaVivo}{\faIcon{circle}\,\textbf{Intermedio}}}
\newcommand{\dificil}{\textcolor{rojoAlerta}{\faIcon{circle}\,\textbf{Difícil}}}

% ─────────────────────────────────────────────────────────────────────────────
\renewcommand{\iconoBanda}{shield-alt}
\fancyhead[L]{\footnotesize\color{grisTexto}\faIcon{shield-alt}~CYBERSEGURIDAD}
\fancyhead[R]{\footnotesize\color{grisTexto}Guía teórica}
\hypersetup{pdftitle={Ciberseguridad Aplicada},pdfauthor={Prof. César Rodríguez}
  pdfsubject={CYBERSEGURIDAD}}

\begin{document}
\bandaTitulo[fondoOscuro]{Guía teórica --- Semana 02}{Criptografía aplicada y gestión de secretos --- material de estudio}

\begin{center}\small\color{grisTexto}Prof. César Rodríguez \textbar{} ELECTRONICA \textbar{} 1 de noviembre de 2026\end{center}

\begin{tabularx}{\linewidth}{@{}>{\raggedright\arraybackslash}p{4.6cm}X@{}}
\toprule
\textbf{Curso} & CYBERSEGURIDAD --- Ciberseguridad Aplicada \\
\textbf{Semana} & 02 de 10 \\
\textbf{Dedicación} & 4 horas de laboratorio + 2 horas de estudio independiente \\
\textbf{Propósito} & Estudio y comprensión; la evaluación es aparte \\
\bottomrule
\end{tabularx}

\section{\iconotexto{book-open}{Cómo usar esta guía}}

Esta guía es material de estudio para la semana 02: explica por qué se elige cada primitiva, qué propiedad garantiza y qué fallos se repiten al usarla. No es un examen ni sustituye a la práctica: los comandos son ejemplos para entender el mecanismo y la evidencia de la semana se produce en el repositorio del grupo, con las evaluaciones teóricas ya publicadas aparte.

\begin{cajaConcepto}\textbf{No es un examen.} Aquí se explica el porqué de cada decisión. Las evaluaciones teóricas por grupo se emiten por separado; esta guía es el material para prepararlas y para entender la práctica de la semana.
\end{cajaConcepto}

\section{\iconotexto{bullseye}{Objetivo de la semana}}

Aplicar cifrado simétrico y asimétrico, funciones hash y derivación de claves, y diagnosticar una configuración TLS insegura.

\section{\iconotexto{clipboard-list}{Contenidos que se estudian}}

\begin{itemize}[leftmargin=1.4em,itemsep=2pt,topsep=3pt]
  \item Cifrado simétrico AES en modos GCM y CBC, y por qué GCM aporta integridad.
  \item Cifrado asimétrico RSA y curvas elípticas; intercambio de claves.
  \item Hash criptográfico, HMAC y funciones de derivación PBKDF2 y Argon2.
  \item TLS 1.2 y 1.3, cadena de certificados y errores de validación.
  \item Gestión de secretos y detección de claves en repositorios.
\end{itemize}

\section{\iconotexto{brain}{Desarrollo de los temas}}

\needspace{6\baselineskip}

\subsection{1. AES en GCM frente a CBC, y por qué GCM da integridad}

AES es un cifrado simétrico de bloque de 128 bits con claves de 128, 192 o 256 bits. En CBC cada bloque se mezcla con el anterior usando un vector de inicialización aleatorio: eso da confidencialidad, pero un atacante puede modificar un bloque del texto cifrado y provocar un cambio previsible en el texto claro sin que nadie lo detecte. GCM es un modo autenticado: además del texto cifrado produce una etiqueta de autenticación de 128 bits que cubre los datos y el contenido adicional. Si un solo byte cambia, la etiqueta no cuadra y el descifrado falla.

\textbf{Ideas clave}

\begin{itemize}[leftmargin=1.4em,itemsep=2pt,topsep=3pt]
  \item CBC necesita relleno, IV aleatorio e imprevisible, y una MAC separada (encrypt-then-MAC) si se quiere integridad.
  \item GCM necesita un nonce único por clave: repetirlo con la misma clave destruye la seguridad de la etiqueta.
  \item El contenido adicional se autentica pero no se cifra; sirve para proteger cabeceras que deben verse.
  \item Confidencialidad e integridad son propiedades distintas: cifrar no significa detectar manipulación.
\end{itemize}

\textbf{Comandos de ejemplo}

\begin{lstlisting}[style=estiloCodigo]
python3 - <<'PY'
from cryptography.hazmat.primitives.ciphers.aead import AESGCM
import os
clave = AESGCM.generate_key(bit_length=256)
nonce = os.urandom(12)
cifrado = AESGCM(clave).encrypt(nonce, b'mensaje de prueba', None)
print('cifrado:', cifrado.hex())
alterado = bytearray(cifrado)
alterado[0] ^= 1
try:
    AESGCM(clave).decrypt(nonce, bytes(alterado), None)
    print('ERROR: se descifro sin detectar el cambio')
except Exception as exc:
    print('integridad rota:', type(exc).__name__)
PY
\end{lstlisting}

\textbf{Errores frecuentes}

\begin{itemize}[leftmargin=1.4em,itemsep=2pt,topsep=3pt]
  \item Reutilizar nonce y clave en GCM, aunque el mensaje sea distinto.
  \item Usar CBC sin MAC y creer que el relleno válido demuestra integridad.
  \item Guardar clave, nonce y texto cifrado en el mismo archivo sin anotar qué es cada cosa.
  \item Probar la integridad con el texto cifrado original en vez de con una copia alterada.
\end{itemize}

\needspace{6\baselineskip}

\subsection{2. RSA, curvas elípticas e intercambio de claves}

RSA y las curvas elípticas son criptografía asimétrica: una clave pública que se distribuye y una privada que nunca sale. RSA apoya su seguridad en la dificultad de factorizar; una curva como P-256 ofrece seguridad comparable con claves mucho más cortas. En la práctica RSA se usa sobre todo para firmar y autenticar, con relleno PSS para firmas u OAEP para cifrado; nunca se firma con RSA crudo. El intercambio de claves lo hace ECDH: cada parte aporta un valor efímero y ambas llegan al mismo secreto sin enviarlo. Como los valores son efímeros, comprometer la clave privada de largo plazo no revela sesiones pasadas: eso es secreto perfecto hacia adelante.

\textbf{Ideas clave}

\begin{itemize}[leftmargin=1.4em,itemsep=2pt,topsep=3pt]
  \item La firma demuestra origen e integridad del mensaje; no lo cifra.
  \item RSA-2048 es el mínimo aceptable hoy; P-256 es el equivalente habitual en curvas.
  \item ECDHE ephemeral es el mecanismo de acuerdo de claves; el certificado solo autentica a quién corresponde esa clave.
  \item Verificar con la clave pública correcta es tan importante como firmar con la privada correcta.
\end{itemize}

\textbf{Comandos de ejemplo}

\begin{lstlisting}[style=estiloCodigo]
openssl genpkey -algorithm RSA -pkeyopt rsa_keygen_bits:2048 -out firma.pem
openssl pkey -in firma.pem -pubout -out firma.pub.pem
openssl dgst -sha256 -sign firma.pem -out mensaje.sig mensaje.txt
openssl dgst -sha256 -verify firma.pub.pem -signature mensaje.sig mensaje.txt
openssl genpkey -algorithm EC -pkeyopt ec_paramgen_curve:P-256 -out curva.pem
\end{lstlisting}

\textbf{Errores frecuentes}

\begin{itemize}[leftmargin=1.4em,itemsep=2pt,topsep=3pt]
  \item Firmar con RSA sin relleno PSS y cifrar con RSA sin OAEP.
  \item Confundir la clave pública del certificado con la clave privada de la sesión.
  \item Reutilizar los mismos parámetros efímeros en varias sesiones.
  \item Creer que una firma válida significa que el mensaje era secreto.
\end{itemize}

\needspace{6\baselineskip}

\subsection{3. Hash, HMAC, PBKDF2 y Argon2}

Un hash criptográfico resume un mensaje en una huella de tamaño fijo: sirve para detectar cambios, no para ocultar el contenido. Un HMAC añade una clave secreta al hash y permite demostrar que el mensaje fue autenticado por quien conoce esa clave. Para contraseñas no basta un hash rápido: hace falta una función de derivación lenta con sal única, porque el atacante prueba millones de candidatos offline. PBKDF2 aplica muchas iteraciones de HMAC-SHA256 y su coste se ajusta con el número de iteraciones. Argon2id, además, exige memoria, lo que encarece los intentos masivos con GPU. Se guarda el algoritmo, los parámetros, la sal y el resultado; nunca la contraseña.

\textbf{Ideas clave}

\begin{itemize}[leftmargin=1.4em,itemsep=2pt,topsep=3pt]
  \item Hash: detecta cambios accidentales o intencionados, sin secreto.
  \item HMAC: detecta manipulación y falsificación si solo el emisor y el receptor conocen la clave.
  \item PBKDF2: resiste fuerza bruta subiendo iteraciones; mide el tiempo antes de elegir el número.
  \item Argon2id: resiste mejor los ataques paralelos porque fija coste de memoria, tiempo y paralelismo.
\end{itemize}

\textbf{Comandos de ejemplo}

\begin{lstlisting}[style=estiloCodigo]
sha256sum mensaje.txt
python3 - <<'PY'
import hashlib
print(hashlib.pbkdf2_hmac('sha256', b'clave', b'sal-unica', 100000).hex())
PY
pip install argon2-cffi
python3 -c "from argon2.low_level import hash_secret; print(hash_secret(b'clave', b'sal-unica'))"
\end{lstlisting}

\textbf{Errores frecuentes}

\begin{itemize}[leftmargin=1.4em,itemsep=2pt,topsep=3pt]
  \item Usar MD5 o SHA-1 para contraseñas o para firma digital.
  \item Guardar la contraseña en claro junto a la sal y llamar eso hash.
  \item Repetir la misma sal para todos los usuarios.
  \item Elegir iteraciones sin medir: demasiadas rompen el inicio de sesión y demasiadas pocas no defienden.
\end{itemize}

\needspace{6\baselineskip}

\subsection{4. TLS 1.2 frente a TLS 1.3, cadena de certificados y errores de validación}

TLS protege el canal entre cliente y servidor. TLS 1.2 negocia la versión y el algoritmo en varios viajes y admite modos antiguos como CBC y el intercambio de claves con RSA estático. TLS 1.3 recorta la negociación a un viaje, elimina los algoritmos obsoletos y exige cifrado autenticado, de modo que también va cifrada buena parte del handshake. El certificado del servidor se valida siguiendo la cadena: hoja, intermedios y una raíz en la tienda de confianza. No basta con que el certificado exista: se comprueba firma, fechas, nombre del servidor en el SAN, uso previsto y revocación.

\textbf{Ideas clave}

\begin{itemize}[leftmargin=1.4em,itemsep=2pt,topsep=3pt]
  \item TLS 1.3 usa nombres de cifrado distintos y solo modos AEAD como AES-GCM o ChaCha20-Poly1305.
  \item La cadena incompleta es un fallo del servidor: el cliente no puede reconstruir los intermedios por arte de magia.
  \item Un certificado autofirmado puede ser válido en el laboratorio si se importa explícitamente, nunca por defecto.
  \item 0-RTT de TLS 1.3 ahorra un viaje pero puede repetirse: no conviene para acciones no idempotentes.
\end{itemize}

\textbf{Comandos de ejemplo}

\begin{lstlisting}[style=estiloCodigo]
openssl s_client -connect example.com:443 -servername example.com -showcerts </dev/null 2>/dev/null | openssl x509 -noout -subject -issuer -dates
openssl x509 -in cert.pem -noout -text | grep -E 'Signature Algorithm|Public-Key'
openssl verify -CAfile ca.pem cert.pem
\end{lstlisting}

\textbf{Errores frecuentes}

\begin{itemize}[leftmargin=1.4em,itemsep=2pt,topsep=3pt]
  \item Confiar en el common name cuando el navegador exige el SAN.
  \item Ignorar un error de cadena con -k o desactivar la verificación para que la prueba pase.
  \item Suponer que un certificado caducado se arregla cambiando la hora del cliente.
  \item Creer que TLS 1.3 permite los mismos cifrados que TLS 1.2.
\end{itemize}

\needspace{6\baselineskip}

\subsection{5. Gestión de secretos y detección de claves en repositorios}

Un secreto es cualquier dato que permite actuar en nombre del sistema: tokens de API, contraseñas, claves privadas o cookies de sesión. No pertenece al código fuente ni al historial de Git: se entrega por variables de entorno, un gestor de secretos o un archivo local ignorado. Si un secreto llegó a un commit, lo primero es rotarlo en el proveedor, porque borrar el commit no lo convierte en secreto otra vez; después se reescribe el historial y se avisa a quienes clonaron. La detección se hace con patrones y entropía antes del commit y también sobre todo el historial, incluidas ramas y etiquetas.

\textbf{Ideas clave}

\begin{itemize}[leftmargin=1.4em,itemsep=2pt,topsep=3pt]
  \item Un .env ignorado protege el futuro, no el pasado: revisa también commits antiguos.
  \item Rotar primero, limpiar después: el orden inverso deja una ventana de abuso.
  \item Los secretos de prueba también se rotan: los bots los encuentran igual.
  \item La huella de un repositorio público incluye los archivos borrados en el historial.
\end{itemize}

\textbf{Comandos de ejemplo}

\begin{lstlisting}[style=estiloCodigo]
git grep -nE '(AKIA[0-9A-Z]{16}|ghp_[A-Za-z0-9]{36}|-----BEGIN [A-Z ]*PRIVATE KEY-----)' || true
git log --all -p -S'BEGIN PRIVATE KEY' -- .
gitleaks detect --source . --redact
\end{lstlisting}

\textbf{Errores frecuentes}

\begin{itemize}[leftmargin=1.4em,itemsep=2pt,topsep=3pt]
  \item Borrar el archivo del último commit y dar el caso por cerrado.
  \item Poner el token en una variable de entorno pero dejar el valor por defecto en el código.
  \item Subir capturas de pantalla o logs con el token visible.
  \item Confiar solo en un escáner: las claves con formatos nuevos se escapan.
\end{itemize}

\needspace{8\baselineskip}

\section{\iconotexto{table}{Mapa rápido}}

\begin{center}
\begin{minipage}{\linewidth}
\centering
{\footnotesize
\textbf{Mapa rápido de primitivas de la semana}\par\smallskip
\begin{tabular}{@{}>{\raggedright\arraybackslash}p{2.40cm}>{\raggedright\arraybackslash}p{3.00cm}>{\raggedright\arraybackslash}p{4.80cm}>{\raggedright\arraybackslash}p{4.40cm}@{}}
\toprule
\textbf{Primitiva} & \textbf{Propiedad principal} & \textbf{No cubre} & \textbf{Uso en la práctica} \\
\midrule
AES-256-GCM & Confidencialidad e integridad & No distribuye la clave entre desconocidos & Cifrar y alterar el archivo \\
Firma RSA o ECC & Origen e integridad del mensaje & No oculta el contenido & Firmar y verificar con el caso alterado \\
HMAC & Autenticación con clave compartida & No sirve para derivar claves de contraseñas & Comparar con un hash sin clave \\
PBKDF2 y Argon2id & Resistencia a fuerza bruta & No cifra ni firma & Medir coste e iteraciones \\
TLS 1.3 & Canal cifrado y autenticado & No arregla un certificado inválido & Leer emisor, validez y algoritmo \\
\bottomrule
\end{tabular}
}
\end{minipage}
\end{center}

\section{\iconotexto{clipboard-check}{Checklist antes de la práctica}}

\begin{itemize}[leftmargin=1.4em,itemsep=2pt,topsep=3pt]
  \item Puedo explicar con mis palabras por qué GCM detecta un byte alterado y CBC no, sin MAC.
  \item Sé distinguir firma, HMAC y hash por la propiedad que aporta y por el secreto que usa.
  \item He medido el tiempo de PBKDF2 con varias iteraciones y puedo justificar el valor elegido.
  \item Sé reconocer en una salida de openssl el emisor, la validez y el algoritmo del certificado.
  \item Compruebo la cadena completa y el SAN antes de aceptar un certificado del laboratorio.
  \item Si aparece un secreto en Git, rotarlo primero y reescribir el historial después.
\end{itemize}

\clearpage

\section{\iconotexto{book}{Glosario de la semana}}

\small
\begin{longtable}{@{}>{\raggedright\arraybackslash}p{4.2cm}>{\raggedright\arraybackslash}p{\dimexpr\linewidth-4.2cm-12pt\relax}@{}}
\toprule
\textbf{Término} & \textbf{Definición}\\
\midrule
\endhead
\bottomrule
\endlastfoot
\textbf{Cifrado simétrico} & Cifrado que usa la misma clave para cifrar y descifrar. Es rápido y adecuado para volumen, pero exige repartir la clave por un canal seguro.\\
\textbf{Cifrado asimétrico} & Cifrado con un par de claves: la pública cifra o verifica y la privada descifra o firma. Resuelve el reparto de claves a cambio de ser más lento.\\
\textbf{Hash criptográfico} & Función de un solo sentido que produce una huella de tamaño fijo; cambiar un byte cambia la huella. No cifra ni se puede revertir.\\
\textbf{HMAC} & Hash con clave que autentica un mensaje y detecta manipulación. Aporta integridad y autenticidad de origen, no confidencialidad.\\
\textbf{Derivación de claves} & Convertir una contraseña o secreto débil en una clave del tamaño adecuado con un coste deliberado (PBKDF2, Argon2) para frenar la fuerza bruta.\\
\textbf{TLS} & Protocolo que cifra y autentica una conexión entre cliente y servidor. HTTPS es HTTP sobre TLS.\\
\textbf{Cadena de certificados} & Secuencia de certificados firmados que enlaza la clave del servidor con una autoridad de confianza. Si falta un eslabón o es inválido, la validación falla.\\
\textbf{Gestión de secretos} & Crear, almacenar, rotar y revocar claves y contraseñas sin dejarlas en el código, en los argumentos ni en el historial de git.\\
\end{longtable}

\begin{cajaRecuerda}La práctica de esta semana pide demostrar cada propiedad con una salida real: cifrar, alterar, firmar, verificar, inspeccionar y medir. Esta guía da el vocabulario; el repositorio del grupo da la evidencia.
\end{cajaRecuerda}
\end{document}
